\chapter{Magnetism and Forces on Charges}

A compass needle swings to point north. A fridge magnet holds up your
report card. An MRI machine looks inside your body without a single cut.
And every electric motor, from a phone's vibrator to a Tesla, spins
because of the same force. That force is magnetism, and it turns out to be
electricity in motion.

\section{Magnets}\index{magnet}
Every magnet has two ends, called its \newterm{north pole} and
\newterm{south pole}. Like poles repel and opposite poles attract, just as
with charges. But there is one big difference: you can never get a north
pole by itself. Cut a bar magnet in half, and each half grows a new pair
of poles.

Earth itself is a giant magnet. A compass needle is a tiny magnet free to
turn, and its north pole points toward Earth's geographic north. (So, oddly,
the magnetic pole near Earth's North Pole is actually a magnetic
\emph{south} pole: it attracts the north end of the needle.)

Where does a magnet's magnetism come from? Every electron acts like a
tiny magnet. In most materials, these tiny magnets point every which way
and cancel out. In iron, nickel, and cobalt, neighboring atoms line up in
regions called \newterm{domains}. When most of the domains point the same
way, the whole piece of metal becomes a magnet.

\section{Magnetic Fields}\index{magnetic field}
A \newterm{magnetic field} $\vec{B}$ is a vector field that exerts forces on moving charges and on
currents. Its unit is the \newterm{tesla} (T). Magnets and currents both produce magnetic fields.
The electric field of the charge chapter pushes on a charge whether or not it moves. A magnetic field
pushes only on a charge that is moving, and the push is perpendicular to both the velocity and the field.

\section{The Field of a Long Wire}\index{Biot-Savart law}
A long straight wire carrying a current $I$ produces a magnetic field that circles the wire. At distance $r$ from the wire,
its magnitude is
$$B = \frac{\mu_0 I}{2\pi r}, \qquad \mu_0 = 4\pi \times 10^{-7}\ \text{T}\cdot\text{m/A}$$
The direction is given by a right-hand rule: point your right thumb along
the current, and your fingers curl in the direction of $\vec{B}$. Hans
Christian {\O}rsted discovered this field in 1820 when he noticed a compass
needle twitch every time he switched on a nearby current. It was the first
clue that electricity and magnetism are connected.

\subsection{A Worked Example}
A wire carries $I = 10$ A. At $r = 0.05$ m,
$$B = \frac{(4\pi \times 10^{-7})(10)}{2\pi(0.05)} = 4 \times 10^{-5} \text{ T}$$
For comparison, Earth's magnetic field at its surface is roughly $5 \times 10^{-5}$ T, so the field of this wire at 5 cm is about the same size as Earth's.

\section{Force on a Moving Charge}\index{Lorentz force}
A charge $q$ moving with velocity $\vec{v}$ in a magnetic field feels the \newterm{magnetic force}
$$\vec{F} = q\vec{v} \times \vec{B}, \qquad |\vec{F}| = |q|vB\sin\theta$$
where $\theta$ is the angle between $\vec{v}$ and $\vec{B}$. The force is zero when the charge moves parallel to the field,
and largest when the motion is perpendicular to it. The cross product means the force is perpendicular to both $\vec{v}$ and $\vec{B}$,
so the magnetic force changes the direction of motion but does no work and does not change the speed.

To find the direction for a positive charge, use another right-hand rule:
point your fingers along $\vec{v}$, curl them toward $\vec{B}$, and your
thumb points along $\vec{F}$. For a negative charge, such as an electron,
the force points the opposite way.

\subsection{A Worked Example}
A proton ($q = 1.6 \times 10^{-19}$ C) moves at $3 \times 10^5$ m/s perpendicular to a 0.5 T field. The force has magnitude
$$F = (1.6 \times 10^{-19})(3 \times 10^5)(0.5) = 2.4 \times 10^{-14} \text{ N}$$

\section{Circular Motion}
If a charge moves perpendicular to a uniform field, the magnetic force always points toward the center of a circle, and the charge
moves in a circle of radius
$$r = \frac{mv}{|q|B}$$
An electron ($m = 9.11 \times 10^{-31}$ kg) moving at $2 \times 10^6$ m/s in a 0.01 T field has a radius of about $1.1 \times 10^{-3}$ m.
Its frequency of revolution is $f = \frac{|q|B}{2\pi m} \approx 2.8 \times 10^8$ Hz, which does not depend on its speed.

Where does the radius formula come from? From the circular motion you
know: the magnetic force $|q|vB$ supplies the centripetal force
$\frac{mv^2}{r}$. Setting them equal and solving for $r$ gives
$r = \frac{mv}{|q|B}$. This is how the auroras form: charged particles from
the Sun spiral along Earth's magnetic field lines down toward the poles,
where they crash into the air and make it glow.

\section{Force on a Current-Carrying Wire}
A wire of length $L$ carrying current $I$ in a field $B$ feels a force
$$F = ILB\sin\theta$$
where $\theta$ is the angle between the wire and the field. This follows from the force on each moving charge in the wire, since the current is
the total charge per second moving through it. A 5 A current in a 0.2 m wire in a 0.3 T field, perpendicular to the field, feels a force of $0.3$ N.

\section{Solenoids}
Wrap a wire into a coil of $n$ turns per meter, and a current $I$ gives a nearly uniform field inside the coil:
$$B = \mu_0 n I$$
For $n = 1000$ turns/m and $I = 2$ A, $B \approx 2.5 \times 10^{-3}$ T. A solenoid is an \newterm{electromagnet}: a magnet you can switch on and
off. Electromagnets lift scrap cars in junkyards, lock doors, and ring
doorbells. The giant superconducting solenoid in an MRI machine makes a
uniform field of 1.5 to 3 T, about 50,000 times stronger than Earth's.

\section{What Comes Next}
A changing magnetic field induces a current in a nearby conductor. That effect, and its direction (Lenz's law), is covered in the chapter on electromagnetic induction.

\begin{Exercise}[label=mag1, title=Field of a wire]
A wire carries $I = 3$ A. Find the magnetic field at $r = 0.02$ m.
\end{Exercise}
\begin{Answer}[ref=mag1]
$B = \frac{(4\pi \times 10^{-7})(3)}{2\pi(0.02)} = 3 \times 10^{-5}$ T.
\end{Answer}

\begin{Exercise}[label=mag2, title=Force on a charge]
A charge $q = 2 \times 10^{-6}$ C moves at 4 m/s perpendicular to a 0.25 T field. Find the magnitude of the magnetic force.
\end{Exercise}
\begin{Answer}[ref=mag2]
$F = qvB\sin 90^\circ = (2 \times 10^{-6})(4)(0.25) = 2 \times 10^{-6}$ N.
\end{Answer}

\begin{Exercise}[label=mag3, title=Force on a wire]
A wire of length 0.5 m carries 2 A and sits at 30 degrees to a 0.4 T field. Find the force on the wire.
\end{Exercise}
\begin{Answer}[ref=mag3]
$F = (2)(0.5)(0.4)\sin 30^\circ = 0.2$ N.
\end{Answer}

\begin{Exercise}[label=mag4, title=Does the field do work?]
A charged particle moves in a circle in a uniform magnetic field. Does the magnetic force change its kinetic energy? Explain.
\end{Exercise}
\begin{Answer}[ref=mag4]
No. The magnetic force is always perpendicular to the velocity, so it does no work. The speed, and therefore the kinetic energy, stays constant.
\end{Answer}

\begin{Exercise}[label=mag5, title=Proton in a circle]
A proton ($m = 1.67 \times 10^{-27}$ kg) moves at $3 \times 10^5$ m/s
perpendicular to a 0.5 T field. What is the radius of its circle? How
does that compare with an electron moving at the same speed?
\end{Exercise}
\begin{Answer}[ref=mag5]
$r = \frac{mv}{|q|B} = \frac{(1.67 \times 10^{-27})(3 \times 10^5)}{(1.6 \times 10^{-19})(0.5)} \approx 6.3 \times 10^{-3}$ m, about 6 mm.
An electron has the same size charge but about $\frac{1}{1836}$ the mass,
so its circle is about 1836 times smaller (about 3.4 micrometers), and it
curves the opposite way.
\end{Answer}
